\subsection*{The rank-129 Ferrari cover and its elliptic Prym blocks}

\paragraph{RP.1 (Proved): character multiplicities and compact dimensions.}
Put $D=C''(C')^2$. Since $C'$ is free of rank five,
$C'/D\simeq(C_2)^5$, $[C:D]=256$, and $D$ is free
of rank $1+32(5-1)=129$. Geometrically, FC.3 identifies
$C'$ with the fundamental group of
$X=E\setminus E[2]$. Its maximal mod-two abelian cover
$Y\to X$ has degree thirty-two. Every puncture inertia
has order two in the deck group $A=(C_2)^5$, so the
compactification of $Y$ has sixty-four punctures removed.
The Euler characteristic is $32(-4)=-128$, giving genus
thirty-three. Thus its compact and boundary homology
dimensions are sixty-six and sixty-three respectively.

Over $\mathbb F_3$, the $A$-action is semisimple. The
trivial character has homology dimension five; each
of the thirty-one nontrivial characters has dimension
four, by the rank-one local-system Euler characteristic
on $X$. If $r_\chi$ is the number of the four puncture
inertias on which $\chi$ is trivial, then its boundary
multiplicity is $r_\chi$ for nontrivial $\chi$ and three
for the trivial character. Indeed the boundary module
is the sum of the four permutation modules induced from
the inertia groups, modulo the one global boundary
relation. Hence the nontrivial compact multiplicity is
$4-r_\chi$.

The four boundary signs have product one. Restriction
to the three-dimensional boundary subspace has a
two-dimensional kernel, coming from elliptic homology.
Consequently there are exactly
\[
\begin{array}{c|c|c|c}
\text{nontrivial character type}&\text{number}&
 \text{proper dimension}&\text{boundary dimension}\\ \hline
\text{unramified at all four punctures}&3&0&4\\
\text{ramified at exactly two punctures}&24&2&2\\
\text{ramified at all four punctures}&4&4&0.
\end{array}
 \tag{RP.1a}
\]
Including the trivial character gives
$2+24\cdot2+4\cdot4=66$ and
$3+3\cdot4+24\cdot2=63$. These are geometric
character counts, not assumptions about the finite
image of the braid action.

\paragraph{RP.2 (Proved): the two-branch covers and their actual Prym curves.}
Let the two selected branch points on $E$ be $Q_1,Q_2$,
both of order two. Translating by $-Q_1$ makes the
branch divisor $O+T$, where $T=Q_2-Q_1\ne O$ is of
order two. The four square roots of this degree-two
line bundle are
\[
                \mathcal O_E(P),\qquad 2P=T.
 \tag{RP.2a}
\]
They give precisely the four double covers with this
fixed branch divisor. Over an algebraically closed
field of characteristic different from two, there is
no additional scalar-twist distinction. A choice of
$P$ is part of the marking and must not be omitted.

Translate once more by $-P$. The branch points become
$-P,P$, and write the translated elliptic curve as
\[
 E:y^2=\prod_{i=1}^3(x-e_i),\qquad a=x(P).
\]
The corresponding genus-two curve has the explicit model
\[
 C_a:y^2=\prod_{i=1}^3(w^2+a-e_i),\qquad x=w^2+a.
 \tag{RP.2b}
\]
The degree-two map to $E$ is the quotient by
$\sigma(w,y)=(-w,y)$. Its two finite fixed points
are exactly the two points over $x=a$. Composing
the quotient with translation by $P$ returns the
original marked branch divisor $O+T$.

The hyperelliptic involution is
$\iota(w,y)=(w,-y)$. The other elliptic quotient,
by $\tau=\iota\sigma$, is
\[
 E'_a:v^2=u\prod_{i=1}^3(u+a-e_i),\qquad
                       u=w^2,\quad v=wy.
 \tag{RP.2c}
\]
It is ramified at the two points at infinity of $C_a$.
Its pullback map $\operatorname{Jac}(E'_a)\to
\operatorname{Jac}(C_a)$ is injective. For a degree-two
ramified map of curves, a nontrivial element of this
kernel would be a line bundle of order two whose
associated connected etale double cover is dominated
by the given double cover; equality of the degrees
would make the original map etale, a contradiction.
The image is anti-invariant under $\sigma$, since
$\tau$ fixes it and $\iota$ acts as minus one on the
Jacobian. It therefore equals the connected Prym
kernel, by dimension. In particular
\[
           \operatorname{Prym}(C_a/E)\simeq E'_a
 \tag{RP.2d}
\]
as elliptic curves. The polarization inherited from
the genus-two Jacobian is twice the principal
elliptic polarization; it does not change this
identification of the underlying elliptic curve.

\paragraph{RP.3 (Proved): the signed quadratic Legendre transformation.}
Normalize the chosen two-torsion point as $T=(0,0)$ on
\[
       E_\lambda:y^2=x(x-1)(x-\lambda),\qquad
       \lambda\ne0,1.
\]
For a point $P$ with $2P=T$, its $x$-coordinate satisfies
$a^2=\lambda$. This follows directly from the tangent
line through $P$ and $T$: substituting a line $y=mx$
into the cubic shows that its nonzero intersection
roots have product $\lambda$, and tangency makes both
equal to $a$. Conversely the two points
\[
        P=(a,\pm i\,a(a-1)),\qquad a^2=\lambda,
 \tag{RP.3a}
\]
have $2P=T$. Thus the four choices consist of the two
values of $a$ and the two signs in (RP.3a).

In (RP.2c), put $x=u+a$. Its four branch points on
the $x$-line are $a,0,1,a^2$. The Mobius transformation
and rescaling
\[
       X=\frac{(1-a)x}{x-a},\qquad
       Y=\frac{v(X+a-1)^2}{a(a-1)}
 \tag{RP.3b}
\]
give exactly
\[
       \boxed{\ E'_a:Y^2=a(a-1)X(X-1)(X+a).\ }
 \tag{RP.3c}
\]
For verification, $x=aX/(X+a-1)$ and
\[
 (x-a)x(x-1)(x-a^2)
 =\frac{a^3(a-1)^3 X(X-1)(X+a)}{(X+a-1)^4}.
\]
Thus the new Legendre parameter is $\mu=-a$, together
with the quadratic twist $a(a-1)$. The unsigned formula
$\mu=\pm\sqrt\lambda$ alone loses a nontrivial central
monodromy character.

The complete geometric diagram is over
$S=\mathbb P^1_a\setminus\{0,1,-1,\infty\}$:
the original elliptic homology is pulled back by
$\lambda=a^2$, while the Prym homology is pulled back
by $\mu=-a$ and tensored with the rank-one quadratic
local system of $\sqrt{a(a-1)}$. Its signs around
$0,1,-1,\infty$ are respectively $-1,-1,+1,+1$.
Equivalently, if $T_0,T_1,T_\infty$ are the usual
local Legendre monodromy classes, then the local
Prym classes are
\[
\begin{array}{c|cccc}
 a&0&-1&1&\infty\\ \hline
 H_1(E'_a)&-T_0&T_1&-I&T_\infty.
\end{array}
 \tag{RP.3d}
\]
Here $T_0,T_1$ are nontrivial unipotents and
$T_\infty$ has trace $-2$. The table gives conjugacy
classes of positively oriented local loops. It does
not choose four independent conjugating paths or
identify them with a pre-existing Schreier basis.
For any chosen basepoint and paths, (RP.3b)--(RP.3c)
give an exact simultaneous identification of local
systems, up to the single resulting basis change.

The puncture $a=1$, which the unsigned map $\mu=-a$
would fill, has monodromy $-I$ on the Prym. This is
the precise extra sign retained by the two-branch
blocks. The comparison is not a generic isogeny
between $E_{a^2}$ and $E'_{a}$: at $a=1$ the former
has potentially multiplicative reduction, whereas
the latter is potentially good with the displayed
finite monodromy. An isogeny would preserve that
distinction.

\paragraph{RP.4 (Proved scope; unresolved comparison).}
The geometric formulas above identify the proper
two-dimensional constituents of the rank-129 cover.
They do not by themselves identify a particular
matrix block of the saved braid representation:
that requires its deck character, the choice of
$P$, and one simultaneous path and basis comparison.
Nor do the formulas prove that a chosen character
is fixed by every element of $\widehat{GT}$; the
full characteristic rank-129 cover is the available
functorial object, and its character summands can
be permuted.

The central sign at $a=1$ is genuine marked
boundary data. However RP.8 below proves that
the entire linear rank-129 representation
annihilates every genuine gentle pentagon
defect. Its equations therefore supply no
additional comparison with $\kappa_d$.
No arithmetic Galois comparison is silently
promoted to such a statement here.

\paragraph{Primary geometric input.}
The Ferrari identification with the universal
punctured elliptic curve is Nakamura, Tsunogai,
and Yasuda (2010, Sections 2.9 and 2.15),
\emph{Journal of the Institute of Mathematics
of Jussieu, 9}(2), 431--448,
\texttt{https://doi.org/10.1017/S1474748009000115}.
The genus, character, double-cover, and change-of-variable
calculations in RP.1--RP.3 are proved explicitly above.

\paragraph{RP.5 (Proved): the marked punctures and the intrinsic rank-three extension.}
The open genus-two cover corresponding to a two-branch
character is obtained by removing the preimages of all
four original points of $E[2]$, not only the two branch
points. After the translation used in RP.2, these four
points are the four solutions of $2R=T$: they are
$\pm P$, with $x=a$, and $\pm R$, with $x=-a$.
Thus the six removed points on the model (RP.2b) are
\[
 (0,\pm i a(a-1)),\qquad
 (\varepsilon\sqrt{-2a},\delta i a(a+1))
       \quad(\varepsilon,\delta\in\{\pm1\}).
 \tag{RP.5a}
\]
All these points are distinct over
$S=\mathbb P^1\setminus\{0,1,-1,\infty\}$.
Write $C_a^\circ$ for this punctured curve.

The four-dimensional deck-character block is
$H_1(C_a^\circ,\mathbb F_3)^-$ for
$\sigma(w,y)=(-w,y)$. Let $d_+$ and $d_-$ be the
differences of small positively oriented loops about
the two $w$-signs in (RP.5a), with the $y$-sign fixed.
They form its two-dimensional boundary subspace.
The hyperelliptic involution $\iota(w,y)=(w,-y)$
interchanges $d_+$ and $d_-$. It acts as minus one
on all proper genus-two homology. Consequently
\[
 H_1(C_a^\circ,\mathbb F_3)^-
       =L_{\rm bdry}\oplus V_3,
 \tag{RP.5b}
\]
where $L_{\rm bdry}$ is the $+1$ eigenspace of
$\iota$, generated by $d_++d_-$, and $V_3$ is its
$-1$ eigenspace. The latter has proper dimension
two and boundary dimension one, the boundary line
being generated by $d_+-d_-$. The hyperelliptic
involution is canonical for a smooth genus-two
curve; thus this splitting is respected by
parallel transport in the stabilizer of the
chosen double-cover character.

Under the quotient by $\tau=\iota\sigma$, the
removed points map to three points on $E'_a$:
the chosen origin $O$ and a pair $Q_a,-Q_a$.
Using the Legendre coordinates (RP.3b),
\[
 X(Q_a)=\frac{1-a}{2},\qquad
 Y(Q_a)=\frac{a^2-1}{4}\sqrt{2a}.
 \tag{RP.5c}
\]
The other choice of the square root gives $-Q_a$.
The involution induced by $\sigma$ is $[-1]$ on
$E'_a$. Pullback under the degree-two quotient,
whose degree is invertible modulo three, identifies
\[
 V_3\simeq H_1(E'_a\setminus\{O,Q_a,-Q_a\},
                         \mathbb F_3)^{[-1]=-1}.
 \tag{RP.5d}
\]
This remains true though the quotient ramifies at
two interior points: homology of a finite quotient
in coefficients of characteristic prime to the
group order is the invariant homology. On the
boundary, pullback sends the difference of the
$Q_a,-Q_a$ loops to a nonzero scalar multiple of
$d_+-d_-$.

The boundary line in $V_3$ is the quadratic local
system $\chi_{2a}$: a loop exchanging the two
square roots in (RP.5c) exchanges $Q_a,-Q_a$ and
negates their boundary difference. Its signs
around $0,1,-1,\infty$ are $-1,+1,+1,-1$.
Tensoring the whole extension by $\chi_{2a}$
makes this boundary line constant. The corresponding
proper elliptic curve and marked rational section are
\[
 \boxed{\quad
 E''_a:Z^2=\frac{a-1}{2}X(X-1)(X+a),\qquad
 Q''_a=\left(\frac{1-a}{2},\frac{a^2-1}{4}\right).
 \quad}
 \tag{RP.5e}
\]
Over the quadratic cover adjoining $\sqrt{2a}$,
this is the change $Z=Y/\sqrt{2a}$. Thus it is an
exact quadratic-twist identification, not a choice
of unrelated isomorphic fibres.

Dually, the boundary/proper extension is the
relative homology extension of an arc from
$-Q''_a$ to $Q''_a$. If a path and homology basis
are chosen at one basepoint, the difference
between an arc and its continuation is an
integral elliptic cycle; these differences give
the extension cocycle. Its Abel--Jacobi section
is $Q''_a-(-Q''_a)=2Q''_a$. This is the precise
marked-point extension meant here by its Kummer
class. Neither its splitting nor the value of
its reduction modulo three is inferred merely
from the proper elliptic representation.

On the original elliptic base, the hyperelliptic
involution in this construction is $R\mapsto T-R$.
It is therefore the lift of the involution
$[-1]+T$ that fixes the chosen square-root line
bundle $\mathcal O(P)$. In the free-product
description this is one of the three marked
involutions $t_i$; multiplying the lift by the
double-cover deck involution exchanges its
$1$- and $3$-dimensional eigenspace labels but
does not change the splitting. The independent
finite certificate
\texttt{check\_rank129\_open\_prym.py} verifies
this $1+3$ decomposition for each of the twenty-four
two-branch characters and verifies that the
$a,b,c$ actions preserve it between characters.
An explicit choice of character and paths is
still needed to equate an individual matrix
entry with the cocycle of $2Q''_a$.

\paragraph{RP.6 (Proved): the section is generically non-torsion.}
Equation (RP.5e) has the integral Weierstrass model
\[
 \mathcal E_a:
 y^2=x^3+2(a-1)^2x^2-4a(a-1)^2x,
 \qquad
 Q_a^{\rm int}=(-(a-1)^2,(a-1)^2(a+1)),
 \tag{RP.6a}
\]
obtained by $x=2(a-1)X$, $y=4(a-1)Z$.
At $a=2$ this specializes to
\[
 E:y^2=x^3+2x^2-8x,\qquad Q=(-1,3),\qquad
                    2Q=(9/4,15/8).
\]
Its discriminant is $2^{12}3^2$, so five and seven
are good primes. Direct enumeration and the group
law give
\[
\begin{array}{c|cc}
 p&\#E(\mathbb F_p)&\operatorname{ord}(\bar Q)\\ \hline
 5&8&4\\
 7&12&6.
\end{array}
 \tag{RP.6b}
\]
If $Q$ were torsion, its order would be divisible
by three from its reduction at seven. Its
three-primary part injects at the good prime
five, a contradiction because that reduced
elliptic curve has order eight. Prime-to-$p$
torsion injectivity here follows directly from
the formal-group filtration of the good
reduction kernel, whose finite successive
quotients are $p$-groups.

Thus $Q$ and $2Q$ are non-torsion. A generic
torsion section would remain torsion at this
smooth specialization, so (RP.5e) is generically
non-torsion as well. The exact arithmetic is
recorded independently in
\texttt{check\_prym\_section.py} and its output.
This argument does not assert that the section
is geometrically indivisible by three; that
would be a separate monodromy or height claim.

\paragraph{RP.7 (Proved): exact CM specializations and their comparison scope.}
The new proper elliptic family has invariant
\[
             j(E''_a)=
       \frac{256(a^2+a+1)^3}{a^2(a+1)^2}=j(-a).
\]
If $a^2+a+1=0$, then (RP.6a), after $X=x-2a$,
becomes
\[
                 y^2=X^3+8,\qquad Q=(a,3).
 \tag{RP.7a}
\]
The point is the order-three CM translate of
$(1,3)$. If $a=-2$, after $X=x+6$ the family becomes
\[
                 y^2=X^3-36X,\qquad Q=(-3,-9).
 \tag{RP.7b}
\]
Both displayed points are non-torsion. For the
first curve, the reduction of $(1,3)$ has order
six at five while the whole reduced curve has
order sixteen at thirteen. For the second,
the point has order six at eleven while the
whole reduced curve has order eight at five.
The same prime-to-$p$ argument proves both
claims; the certificate in RP.6 checks them.

Relative to the literal R4 models
$E_3:y^2=x^3+1$ and $E_4:y^2=x^3-x$,
(RP.7a) is the quadratic twist by two and
(RP.7b) is the quadratic twist by six.
The exact isomorphisms over the corresponding
quadratic fields are
\[
 (X,y)\longmapsto\left(\frac X2,
                            \frac y{2\sqrt2}\right)
          \quad\hbox{to }E_3,
 \qquad
 (X,y)\longmapsto\left(\frac X6,
                            \frac y{6\sqrt6}\right)
          \quad\hbox{to }E_4.
 \tag{RP.7c}
\]
They carry the two marked points to
\[
 \left(\frac a2,\frac3{2\sqrt2}\right)\in E_3,
 \qquad
 \left(-\frac12,-\frac3{2\sqrt6}\right)\in E_4.
\]
The Tate-module representations before these
isomorphisms differ from those of the R4
curves by the quadratic characters of two
and six. Both characters restrict trivially
to $G_{\mathbb Q^{\rm ab}}$, since that group
fixes the quadratic fields. Thus the proper
torsion representations agree there after
the explicit identifications (RP.7c).
The displayed marked points remain
non-torsion; their division-point extensions
and path cocycles are additional data.
Triviality of the quadratic twists on the
arithmetic subgroup is not a construction
of a common marked path for all of
$\widehat{GT}$.

The corresponding original base curves
$E_{a^2}$ have respectively $j=2048/3$ and
$j=35152/9$. In particular the two elliptic
slots in this construction must be kept
distinct. The CM property of the Prym at
these values does not replace the non-torsion
marked section by the old torsion-frame data.

The characteristic rank-129 cover, its
Prym section, and its finite braid action
are genuine. A comparison of its section's
path cocycle with $\kappa_d$ would still
require a marked full-$\widehat{GT}$ theorem;
the arithmetic specialization does not
provide one. Moreover RP.8 proves that
the linear pentagon in this representation
is automatic on the gentle group. Thus
it cannot itself supply the missing theorem,
even though its harmonic pair image has
non-nilpotent second derived subgroup.

\paragraph{RP.8 (Proved): the entire linear cover is blind to the gentle pentagon.}
The rank-129 construction passes the pair-second-derived
nilpotence test but ultimately contributes no new pentagon
equation. The following source identities explain why.
Write $u=\xi_1$, $v=\xi_2$, $s=aba$, and
$X_4=s^2=z_3$. The exact Ferrari Nielsen formulas give
\[
 c=ua,\qquad a(u)=u,\qquad
 s(u)=v^{-1},\qquad s(v)=vuv^{-1},\qquad
 cbc=uv^{-1}s.
\]
Consequently
\[
 X_3X_1^{-1}=u^2\in C',\qquad
 X_5X_4^{-1}=uv^{-1}u^{-1}v^{-1}\in C'.
 \tag{RP.8a}
\]
The second membership also follows from its exponent sums
$(0,-2)$ and $C'=K'K^2$. Since $X_4$ is central in
the actual three-strand section, it commutes with
$X_1,X_2$. Thus in $L/C'$ the five inputs are
$(X_1,X_2,X_1,X_4,X_4)$, with $X_4$ central among
them. Every antisymmetric commutator word therefore
has pentagon one in this quotient, and
\[
                         \mathcal P(f)\in C'.
 \tag{RP.8b}
\]

Consider the twenty-four rank-three blocks of RP.5.
On the block for a character $\chi:C'/D\to\{\pm1\}$,
every $g\in C'$ acts as the scalar $\chi(g)$, both
on the full block and its boundary line. Tensoring
the block with the dual of that line cancels this
scalar. On the direct sum, with its genuine
permutation of the characters, this is a representation
of $L$ which kills $C'$, hence factors through $L/C'$.
Its pentagon is therefore identically one by (RP.8a).

The line normalization does not preserve an arbitrary
commutator-word value; this qualification is necessary.
It does preserve every actual multiplier-one gentle
$f$. Indeed such an $f$ belongs to $\gamma_3(F)$:
in the class-two free quotient write $f=[x,y]^e$;
the three-cycle relation gives $3e=0$ in
$\widehat{\mathbb Z}$, and hence $e=0$.
The exact signed boundary-line pair images on all
twenty-four characters have nilpotency class at most
two. The two harmonic pair images have lower-central
orders $8,2,1$; the other three have orders $16,2,1$.
Thus every actual $f$ has boundary-line value one
in all five factors. The normalized and original
rank-three pentagon values coincide.

These boundary assertions are finite signed-permutation
identities, not estimates. They are checked on the
whole twenty-four-line representation, including all
its pair orbits. Reduction modulo three is faithful
on signed permutation matrices, so the same identities
hold in the integral boundary lattice and after any
odd coefficient change. An independent audit also
checks directly that each pair commutator commutes
with both generators on all twenty-four lines.

It follows that $\mathcal P(f)$ acts trivially on
each of the twenty-four rank-three blocks. By
(RP.8b), that action is the scalar
$\chi(\mathcal P(f))$. The twenty-four two-branch
characters span $(C'/D)^*$: differences of two
characters with the same boundary restriction
give the two unramified handle characters, and
the six weight-two boundary patterns span the
three-dimensional boundary dual. Therefore
\[
                \boxed{\ \mathcal P(f)\in D
                    =C''(C')^2\quad
                    \text{for every genuine gentle }(1,f).\ }
 \tag{RP.8c}
\]
Every element of $D$ acts trivially by conjugation
on $D^{\rm ab}$. Hence the pentagon is one on
$H_1(D,R)$ for every coefficient ring $R$, not
only on the tested $\mathbb F_3$ quotient.
The apparent boundary/proper equation in FP.47
is therefore an identity on the actual gentle
locus. The non-nilpotent harmonic image of
FP.45 remains correct, but is insufficient
to make the five-point product detect a new
condition.

This conclusion must be kept separate from
the affine extension
$C/(D'D^3)$. The linear conjugation action
on $H_1(D,\mathbb F_3)$ kills the entire
normal subgroup $D$; the affine extension
retains its class in $D/(D'D^3)$. Equation
(RP.8c) does not kill that class and therefore
does not settle the pentagon in the affine
extension.
